\section*{Capítulo I, Problema XV}

\textbf{Problema:}

15. Sean $A, B, C$ tres vectores no nulos. Si $A \cdot B = A \cdot C$, demuestra con un ejemplo que no necesariamente se tiene $B = C$.

\textbf{Solución:}

Para demostrar que $A \cdot B = A \cdot C$ no implica necesariamente que $B = C$, consideremos el siguiente ejemplo:

Sean los vectores:
\begin{align*}
    A &= (1, 0), \\
    B &= (1, 1), \\
    C &= (1, -1).
\end{align*}

Calculamos los productos escalares $A \cdot B$ y $A \cdot C$:
\begin{align*}
    A \cdot B &= (1)(1) + (0)(1) = 1, \\
    A \cdot C &= (1)(1) + (0)(-1) = 1.
\end{align*}

En este caso, $A \cdot B = A \cdot C$, pero claramente $B \neq C$ ya que $B = (1, 1)$ y $C = (1, -1)$.

Por lo tanto, hemos demostrado que $A \cdot B = A \cdot C$ no implica necesariamente que $B = C$.